\documentclass[11pt]{article}
\usepackage[a4paper,margin=25mm]{geometry}
\usepackage[hidelinks]{hyperref}
\usepackage{microtype}
\setlength{\parindent}{0pt}
\setlength{\parskip}{8pt}
\begin{document}

20 September 2026

Editors-in-Chief\\
\emph{Autonomous Agents and Multi-Agent Systems}

Dear Editors-in-Chief,

Please consider our manuscript, ``Same Words, Different Procedure:
Counterfactual Transfer Tests for Reusable LLM-Agent Skills,'' as a Research
article (Regular Paper).

Reusable-skill evaluations commonly report whether an agent succeeds after a
skill is supplied. That outcome cannot distinguish benefit from base-agent
competence, generic executor assistance, or a misleading surface-similar source.
We introduce the Counterfactual Skill Transfer Test (CSTT), a six-condition
same-state protocol that crosses source procedural compatibility with lexical
similarity and includes no-skill and fair binding-only controls.

Following two retained failed pilots and a final mechanism revision, we froze
the method, source assignment, three local model revisions, statistics,
decision gates, and a 60-target zero-overlap confirmation roster. Across 180
target--model blocks and 1,080 native episodes, a procedurally matched but
lexically far source achieved 97.8\% success, versus 10.0\% for a lexically near
procedural lure. The paired gain was 87.8 percentage points (95\% CI
82.8--93.3; task-clustered randomization $p=1.00\times10^{-6}$). Relative to
the binding-only control, matched skills produced 150 helpful and no harmful
discordances. Three independent audits accepted all episodes and 46,662 native
turns without a state-match, native-action, source-assignment, or
information-boundary failure.

The manuscript is within the journal's scope because it provides an evaluation
method for agent decision-making architectures, an auditable testbed protocol,
and empirical evidence about reliable transfer in autonomous agents. It does
not claim to solve retrieval or establish cross-domain transfer; the evidence
is limited to the released ALFWorld confirmation and three open model families.

For transparency, our related manuscript, ``From Traces to Witnessed Contracts:
Incremental Precondition--Effect Compilation for LLM Agent Skills,'' is under
consideration at \emph{Automated Software Engineering}. It studies
pre-execution detection of defective skill revisions. The present paper studies
causal attribution of transfer benefit through a different protocol. The two
papers have disjoint confirmation targets and different methods, estimands,
experiments, figures, tables, endpoints, and results. We include the related
manuscript for editorial assessment.

The submitted manuscript is original and is not under consideration elsewhere.
It contains no third-party figures or tables requiring permission. The work
received no funding, and the authors have no relevant financial or
non-financial interests to disclose. Generative-AI assistance is described in
the Experimental Design section. The public reproducibility record will include
source code, frozen protocols, native interaction logs, audits, tests, and
analysis outputs; third-party game data and model weights are not redistributed.

Sincerely,

Liang Song\\
Corresponding author\\
Business School, Xi'an International University\\
Xi'an, Shaanxi 710077, China\\
\href{mailto:liangsong_1976@126.com}{liangsong\_1976@126.com}

On behalf of Liang Song and Zhai Jiabao

\end{document}
